\subsection{从基线到差异化模型}

创新先回答“基础方案在哪个可验证环节失败”，再引入复杂度。可执行路径包括：
\begin{enumerate}
  \item \textbf{结构创新：}把题意中的时空、网络、分层或阶段结构显式写入状态和约束；
  \item \textbf{机理—数据融合：}用守恒方程约束参数范围，以数据校准未知项，并分开报告拟合和外推；
  \item \textbf{基线—改进链：}先给解析、贪心、线性或持续性基线，再针对残差、非线性或不确定性改进；
  \item \textbf{约束驱动：}将现实可行性写入模型，并报告活跃约束、不可行诊断和修复代价；
  \item \textbf{不确定性驱动：}把参数误差传播到结果和决策，寻找稳定区间与策略切换点；
  \item \textbf{验证创新：}增加独立回代、网格收敛、消融、滚动验证或反事实情景，使结论比模型名更可信。
\end{enumerate}

不采用“换一个更新颖算法名称”作为独立创新。若鲸鱼、麻雀或蛾火焰优化相对
确定性基线、GA/PSO 没有同预算优势，就不进入主模型；若深度网络没有超过季节
或树模型基线，保留简单方案；若多指标组合权重没有通过扰动和秩稳定检验，不把
组合本身写成贡献。

